\setcounter{section}{1}

  \zwischenueberschrift{Übungsaufgaben}

\inputaufgabegibtloesung
{}
{

Es sei

\mavergleichskette
{\vergleichskette
{ G
}
{ \subseteq }{ \R^n
}
{  }{
}
{    }{
}
{   }{
}
}
{}{}{}
\definitionsverweis {offen}{}{}
und
\maabbdisp {\varphi} { G } { \R^m
} {}
eine
\definitionsverweis {stetig differenzierbare Abbildung}{}{,}
die im Punkt

\mavergleichskette
{\vergleichskette
{ P
}
{ \in }{ G
}
{  }{
}
{    }{
}
{   }{
}
}
{}{}{}
ein surjektives
\definitionsverweis {totales Differential}{}{}
besitze. Es sei

\mavergleichskette
{\vergleichskette
{ v
}
{ \in }{ T_PZ
}
{  }{
}
{    }{
}
{   }{
}
}
{}{}{}
ein Vektor des
\definitionsverweis {Tangentialraumes an die Faser}{}{}
$Z$ zu $\varphi$ durch $P$. Zeige, dass es eine
\definitionsverweis {stetig differenzierbare Kurve}{}{}
\maabbdisp {\gamma} { ]- \epsilon, \epsilon[ } { Z
} {}
\zusatzklammer {für ein geeignetes

\mavergleichskettek
{\vergleichskettek
{ \epsilon
}
{ > }{ 0
}
{  }{
}
{    }{
}
{   }{
}
}
{}{}{}} {} {}
mit

\mavergleichskette
{\vergleichskette
{ \gamma(0)
}
{ = }{ P
}
{  }{
}
{    }{
}
{   }{
}
}
{}{}{}
und mit

\mavergleichskettedisp
{\vergleichskette
{ \gamma'(0)
}
{ =} { v
}
{ } {
}
{ } {
}
{ } {
}
}
{}{}{}
gibt.

}
{} {}

\inputaufgabe
{}
{

Es sei

\mavergleichskette
{\vergleichskette
{ W
}
{ \subseteq }{ \R^n
}
{  }{
}
{    }{
}
{   }{
}
}
{}{}{}
\definitionsverweis {offen}{}{,}
\maabb {h} { W } { \R
} {}
eine
\definitionsverweis {stetig differenzierbare Funktion}{}{}
und

\mavergleichskette
{\vergleichskette
{ Y
}
{ = }{ h^{-1}(0)
}
{  }{
}
{    }{
}
{   }{
}
}
{}{}{}
die
\definitionsverweis {Faser}{}{}
zu $0$, wobei $h$ in jedem Punkt von $Y$
\definitionsverweis {regulär}{}{}
sei. Es sei
\maabbdisp {g} { W } {\R
} {}
eine stetig differenzierbare Funktion, die auf $Y$ keine Nullstelle besitzt. Zeige, dass man $Y$ genauso gut als Faser zu

\mavergleichskettedisp
{\vergleichskette
{ \tilde{h}
}
{ =} { g h
}
{ } {
}
{ } {
}
{ } {
}
}
{}{}{}
erhalten kann, dass die
\definitionsverweis {Gradienten}{}{}
zu $h$ und zu $\tilde{h}$  skalare Vielfache voneinander sind und dass die
\definitionsverweis {Tangentialräume}{}{}
nur von $Y$ abhängen.

}
{} {}

\inputaufgabe
{}
{

Es sei

\mavergleichskette
{\vergleichskette
{ W
}
{ \subseteq }{ \R^n
}
{  }{
}
{    }{
}
{   }{
}
}
{}{}{}
\definitionsverweis {offen}{}{,}
\maabb {h} { W } { \R
} {}
eine
\definitionsverweis {stetig differenzierbare Funktion}{}{}
und

\mavergleichskette
{\vergleichskette
{ Y
}
{ = }{ h^{-1}(c)
}
{  }{
}
{    }{
}
{   }{
}
}
{}{}{}
die
\definitionsverweis {Faser}{}{}
zu

\mavergleichskette
{\vergleichskette
{ c
}
{ \in }{ \R
}
{  }{
}
{    }{
}
{   }{
}
}
{}{}{,}
wobei $h$ in jedem Punkt von $Y$
\definitionsverweis {regulär}{}{}
sei. Es sei
\maabbdisp {L} {\R^n } { \R^n
} {}
eine
\definitionsverweis {lineare Isomorphie}{}{,}

\mavergleichskette
{\vergleichskette
{ W'
}
{ = }{ L^{-1}(W)
}
{  }{
}
{    }{
}
{   }{
}
}
{}{}{}
und

\mavergleichskettedisp
{\vergleichskette
{ Z
}
{ =} { L^{-1}(Y)
}
{ =} { (h \circ L)^{-1}(c)
}
{ } {
}
{ } {
}
}
{}{}{.}
Zeige, dass sich die Konzepte
\definitionsverweis {Tangentialraum}{}{,}
tangentiales Vektorfeld, Lösung zu einer zugehörigen tangentialen Differentialgleichung  auf $Y$ und auf $Z$ im Allgemeinen entsprechen
\zusatzklammer {also mithilfe von $L$ ineinander überführt werden können} {} {.}
Wie sieht es mit dem Gradienten aus?

}
{} {}

\inputaufgabe
{}
{

Es sei

\mavergleichskette
{\vergleichskette
{ Y
}
{ = }{ h^{-1}(c)
}
{  }{
}
{    }{
}
{   }{
}
}
{}{}{}
eine
\definitionsverweis {differenzierbare Hyperfläche}{}{}
und sei
\maabbdisp {f} { U} { \R
} {}
eine auf einer offenen Umgebung

\mavergleichskette
{\vergleichskette
{ Y
}
{ \subseteq }{ U
}
{  }{
}
{    }{
}
{   }{
}
}
{}{}{}
definierte stetig differenzierbare Funktion. Es besitze $f$ in

\mavergleichskette
{\vergleichskette
{ P
}
{ \in }{ Y
}
{  }{
}
{    }{
}
{   }{
}
}
{}{}{}
ein
\definitionsverweis {lokales Extremum}{}{}
auf $Y$. Zeige, dass die
\definitionsverweis {tangentiale Komponente}{}{}
von $\operatorname{Grad} \,  f  (  P )$ gleich $0$ ist.

}
{} {}

\inputaufgabe
{}
{

Es sei

\mavergleichskette
{\vergleichskette
{ Y
}
{ \subseteq }{ \R^3
}
{  }{
}
{    }{
}
{   }{
}
}
{}{}{}
die Einheitskugeloberfläche. Bestimme die
\definitionsverweis {orthogonale Zerlegung}{}{}
des Vektors

\mavergleichskette
{\vergleichskette
{ \begin{pmatrix}  6  \\ 7 \\  -3   \end{pmatrix}
}
{ \in }{ \R^3
}
{  }{
}
{    }{
}
{   }{
}
}
{}{}{}
in die tangentiale und in die orthogonale Komponente zum Punkt

\mavergleichskette
{\vergleichskette
{ \begin{pmatrix}  -  { \frac{   1  }{  2 } }  \\ { \frac{   1  }{  2 } } \\  { \frac{   1  }{  \sqrt{2}  } }   \end{pmatrix}
}
{ \in }{ Y
}
{  }{
}
{    }{
}
{   }{
}
}
{}{}{.}

}
{} {}

\inputaufgabe
{}
{

Es sei

\mavergleichskette
{\vergleichskette
{ Y
}
{ \subseteq }{ \R^3
}
{  }{
}
{    }{
}
{   }{
}
}
{}{}{}
die
\definitionsverweis {Einheitssphäre}{}{.}
Es sei
\maabbdisp {\gamma} {[0, 2 \pi]} { Y
} {}
die gleichförmige Kreisbewegung, die die Sphäre auf der Höhe ${ \frac{   1  }{  2 } }$ durchläuft.
\aufzaehlungzwei {Parametrisiere $\gamma$.
} {Bestimme die
\definitionsverweis {tangentiale Beschleunigung}{}{}
von $\gamma$ zu jedem Zeitpunkt

\mavergleichskette
{\vergleichskette
{ t
}
{ \in }{ [0, 2 \pi]
}
{  }{
}
{    }{
}
{   }{
}
}
{}{}{.}
}

}
{} {}

\inputaufgabe
{}
{

Es sei $Y$ der
\definitionsverweis {Graph}{}{}
zu
\maabbeledisp {f} {\R^2} { \R
} {(x,y)} { x^2+y^3
} {,}
und es sei
\maabbdisp {\gamma} {\R} { Y
} {}
der Weg auf dem Graphen oberhalb der Basiskurve
\mathl{t \mapsto (t,t)}{.} Bestimme zu jedem

\mavergleichskette
{\vergleichskette
{ t
}
{ \in }{ \R
}
{  }{
}
{    }{
}
{   }{
}
}
{}{}{}
die
\definitionsverweis {tangentiale Beschleunigung}{}{}
von $\gamma$ in $\gamma(t)$.

}
{} {}

\inputaufgabe
{}
{

Es sei

\mavergleichskette
{\vergleichskette
{ Y
}
{ \subseteq }{ \R^n
}
{  }{
}
{    }{
}
{   }{
}
}
{}{}{}
eine Hyperebene, also die Faser zu einer linearen Funktion

\mavergleichskettedisp
{\vergleichskette
{ h(x_1 , \ldots , x_n)
}
{ =} { a_1x_1 + \cdots + a_nx_n
}
{ } {
}
{ } {
}
{ } {
}
}
{}{}{}
über

\mavergleichskette
{\vergleichskette
{ c
}
{ \in }{ \R
}
{  }{
}
{    }{
}
{   }{
}
}
{}{}{}
mit

\mavergleichskette
{\vergleichskette
{ a_i
}
{ \neq }{ 0
}
{  }{
}
{    }{
}
{   }{
}
}
{}{}{}
für zumindest einen Index $i$. Interpretiere die Begriffe
\definitionsverweis {totales Differential}{}{,}
\definitionsverweis {Gradient}{}{,}
\definitionsverweis {Tangentialraum}{}{,}
Zerlegung in tangentiale und normale Komponente,
\definitionsverweis {tangentiale Beschleunigung}{}{}
und
\definitionsverweis {geodätische Kurve}{}{}
in diesem Spezialfall.

}
{} {}

\inputaufgabe
{}
{

Es sei

\mavergleichskette
{\vergleichskette
{ Y
}
{ = }{ h^{-1}(c)
}
{  }{
}
{    }{
}
{   }{
}
}
{}{}{}
eine
\definitionsverweis {differenzierbare Hyperfläche}{}{,}
die eine Gerade
\mathl{P+ \R v}{} enthalte. Zeige, dass die gleichförmig durchlaufene Gerade eine
\definitionsverweis {geodätische Kurve}{}{}
auf $Y$ ist.

}
{} {}

Zwei Punkte

\mathbed {P,Q \in K} {}
 {P \neq Q} {}
 {} {} {} {,}
heißen \stichwort {antipodal} {,} wenn ihre Verbindungsgerade durch den Kugelmittelpunkt läuft.

\inputaufgabe
{}
{

Es sei

\mavergleichskette
{\vergleichskette
{ K
}
{ \subseteq }{ \R^3
}
{  }{
}
{    }{
}
{   }{
}
}
{}{}{}
eine Kugeloberfläche. Zeige, dass je zwei nicht
\definitionsverweis {antipodale Punkte}{}{}

\mathbed {P,Q \in K} {}
 {P \neq Q} {}
 {} {} {} {,}
auf genau einem
\definitionsverweis {Großkreis}{}{}
von $K$ liegen.

}
{} {}

\bild{ \begin{center}
\includegraphics[width=5.5cm]{ \bildeinlesungsvg {Planned flight map of the Oiseau Blanc} {svg} }
\end{center}
\bildtext {Warum fliegt das Flugzeug einen Bogen?} }

\bildlizenz { Planned flight map of the Oiseau Blanc.svg } {} {Pethrus} {Commons} {CC-by-sa 3.0} {}

\inputaufgabe
{}
{

Ein Flugzeug soll von Osnabrück aus zu einem Zielort auf der Südhalbkugel fliegen. Kann es kürzer sein, in Richtung Norden zu fliegen?

}
{} {}

\inputaufgabe
{}
{

Lucy Sonnenschein findet, dass die Tage zu kurz sind. Daher entschließt sie sich, jeden Tag, und zwar bei Tageslicht, um eine Zeitzone nach Westen zu fliegen, damit jeder Tag $25$ statt $24$ Stunden besitzt.
\aufzaehlungfuenf{Kann Lucy Sonnenschein mit dieser Strategie den Anteil an Lichtstunden in ihrem Leben erhöhen?
}{Kann Lucy Sonnenschein mit dieser Strategie ihr Leben verlängern?
}{Lucy hat jetzt Freunde in aller Welt. Nach wie vielen Tagen sieht sie sie wieder?
}{Wodurch erlebt Lucy zeitliche Einbußen, die ihren täglichen Zeitgewinn ausgleichen?
}{Hat das was mit Relativitätstheorie zu tun?
}

}
{} {}

\inputaufgabe
{}
{

Man gebe ein Beispiel für eine zweifach differenzierbare Kurve an, die sich auf einem
\definitionsverweis {Großkreis}{}{}
der Einheitssphäre bewegt, die aber keine
\definitionsverweis {geodätische Kurve}{}{}
ist.

}
{} {}

\inputaufgabe
{}
{

Ein \stichwort {geodätisches Dreieck} {} auf der Sphäre

\mavergleichskettedisp
{\vergleichskette
{ S^2
}
{ =} { \left\{ (x,y,z) \in \R^3  \mid   x^2+y^2+z^2 = 1        \right\}
}
{ } {
}
{ } {
}
{ } {
}
}
{}{}{}
besteht aus drei Punkten

\mavergleichskette
{\vergleichskette
{ A,B,C
}
{ \in }{ S^2
}
{  }{
}
{    }{
}
{   }{
}
}
{}{}{,}
wobei je zwei Punkte durch einen Ausschnitt eines
\definitionsverweis {Großkreises}{}{}
verbunden werden
\zusatzklammer {die Punkte seien nicht antipodal} {} {.}
Zeige, dass die Summe der Innenwinkel des Dreiecks größer als $\pi$ ist.

}
{} {}

  \zwischenueberschrift{Aufgaben zum Abgeben}

\inputaufgabe
{3}
{

Es sei $Y$ die Hyperfläche, die durch die Bedingung

\mavergleichskettedisp
{\vergleichskette
{ 2x^2 +3y^2+5z^2
}
{ =} { 10
}
{ } {
}
{ } {
}
{ } {
}
}
{}{}{}
im $\R^3$ gegeben ist. Bestimme die orthogonale Zerlegung des Vektors

\mavergleichskette
{\vergleichskette
{ \begin{pmatrix}  4  \\ 2 \\  -5   \end{pmatrix}
}
{ \in }{ \R^3
}
{  }{
}
{    }{
}
{   }{
}
}
{}{}{}
in die tangentiale und in die orthogonale Komponente zum Punkt

\mavergleichskette
{\vergleichskette
{ \begin{pmatrix}  -1 \\ -1\\  1   \end{pmatrix}
}
{ \in }{ Y
}
{  }{
}
{    }{
}
{   }{
}
}
{}{}{.}

}
{} {}

\inputaufgabe
{4}
{

Es sei $Y$ der
\definitionsverweis {Graph}{}{}
zu
\maabbeledisp {f} {\R^2} { \R
} {(x,y)} { x^2+y^2
} {,}
und es sei
\maabbdisp {\gamma} {\R} { Y
} {}
der Weg auf dem Graphen oberhalb der Basiskurve
\mathl{x \mapsto (x,1)}{.} Bestimme zu jedem

\mavergleichskette
{\vergleichskette
{ x
}
{ \in }{ \R
}
{  }{
}
{    }{
}
{   }{
}
}
{}{}{}
die
\definitionsverweis {tangentiale Beschleunigung}{}{}
von $\gamma$ in $\gamma(x)$.

}
{} {}

\inputaufgabe
{4}
{

Am 11.4.2023 um 17:45 Uhr wird zur allgemeinen Verwunderung plötzlich die Erddrehung um die eigene Achse angehalten. GLeichzeitig werden der Luftwiderstand und sonstige Reibungsverluste abgestellt. Alle Menschen, Tiere und Gegenstände, die nicht fixiert sind oder sich nicht rechtzeitig festhalten, führen ihre momentane Bewegung nach dem Trägheitsgesetz einfach fort. Die Gravitation bleibt bestehen, sodass man zum Glück nicht in das Weltall entweicht. Wann kommt der Tafelschwamm aus dem Osnabrücker Hörsaal
\zusatzklammer {er fliegt aus dem Fenster} {} {}
zum ersten Mal wieder zu seiner Ausgangsposition zurück?

}
{} {}

\inputaufgabe
{4 (3+1)}
{

Es sei

\mavergleichskette
{\vergleichskette
{ W
}
{ \subseteq }{ \R^n
}
{  }{
}
{    }{
}
{   }{
}
}
{}{}{}
\definitionsverweis {offen}{}{,}
\maabb {h} { W } { \R
} {}
eine
\definitionsverweis {stetig differenzierbare Funktion}{}{}
und

\mavergleichskette
{\vergleichskette
{ Y
}
{ = }{ h^{-1}(c)
}
{  }{
}
{    }{
}
{   }{
}
}
{}{}{}
die
\definitionsverweis {Faser}{}{}
zu

\mavergleichskette
{\vergleichskette
{ c
}
{ \in }{ \R
}
{  }{
}
{    }{
}
{   }{
}
}
{}{}{,}
wobei $h$ in jedem Punkt von $Y$
\definitionsverweis {regulär}{}{}
sei. Es sei
\maabbdisp {\gamma} { I } { Y
} {}
eine zweimal
\definitionsverweis {stetig differenzierbare}{}{}
Kurve.
\aufzaehlungzwei {Zeige, dass, wenn $\gamma$ eine
\definitionsverweis {geodätische Kurve}{}{}
ist, dann die
\definitionsverweis {Norm}{}{}
von $\gamma'$ konstant ist,
} {Man gebe ein Beispiel einer Kurve mit konstanter Geschwindigkeitsnorm, die keine geodätische Kurve ist.
}

}
{} {}

\inputaufgabe
{3}
{

Es seien
\maabb {h,g} {\R^n } { \R
} {}
\definitionsverweis {stetig differenzierbare Funktionen}{}{}
und

\mavergleichskette
{\vergleichskette
{ Y
}
{ = }{ h^{-1}(c)
}
{  }{
}
{    }{
}
{   }{
}
}
{}{}{}
bzw.

\mavergleichskette
{\vergleichskette
{ Z
}
{ = }{ h^{-1}(d)
}
{  }{
}
{    }{
}
{   }{
}
}
{}{}{}
die zugehörigen
\definitionsverweis {differenzierbaren Hyperflächen}{}{,}
wobei $h$ in jedem Punkt von $Y$ und $g$ in jedem Punkt von $Z$
\definitionsverweis {regulär}{}{}
sei. Es sei

\mavergleichskette
{\vergleichskette
{ T
}
{ = }{ Y \cap Z
}
{  }{
}
{    }{
}
{   }{
}
}
{}{}{}
und es sei
\maabbdisp {\gamma} { I } { T
} {}
eine zweimal stetig differenzierbare Kurve, die aufgefasst als Kurve in $Y$ eine
\definitionsverweis {geodätische Kurve}{}{}
sei. Ist  $\gamma$ auch eine geodätische Kurve in $Z$?

}
{} {}

 [[Kategorie:Latexseite]]
